\chapter{C++ pour la latence II~: temps de compilation}\label{ch:ll:cpp-for-latency-ii-compile-time}

Quatre étapes d'un chemin critique (hot path), chacune une classe avec une fonction virtuelle appelée à travers un pointeur, coûtent sur l'ordinateur portable de ce livre environ
sept nanosecondes par événement. Les mêmes quatre étapes composées sous forme de paramètres de patron coûtent environ une~: le compilateur, voyant
chaque appel, a inliné les quatre dans treize instructions machine sans aucun appel restant. Le travail était identique~; ce qui
a changé est la quantité d'information dont dispose le compilateur au moment où il compile. Ce chapitre traite du déplacement des décisions du temps d'exécution vers le temps
de compilation en C++~: répartition par patrons au lieu de fonctions virtuelles, valeurs calculées par le compilateur, \omterm{def:ll:cpp-for-latency-ii-compile-time:inline}{inlining} et ce qui l'empêche,
ainsi que la disposition du code rare à l'écart du chemin commun.

\section{Répartition virtuelle et son coût}

\begin{definition}[Répartition dynamique, dévirtualisation]\label{def:ll:cpp-for-latency-ii-compile-time:dyn}
Dans la \emph{répartition dynamique}\index{répartition dynamique} (dynamic dispatch), la fonction appelée dépend du type à l'exécution d'un objet~: un appel C++
virtuel charge la table des pointeurs de fonctions de l'objet et appelle à travers elle, un appel indirect. La \emph{dévirtualisation}\index{dévirtualisation}
est le remplacement par le compilateur d'un tel appel par un appel direct, lorsqu'il peut établir le type de l'objet (une classe \texttt{final},
un objet local, une analyse de programme complet).
\end{definition}

Un appel indirect coûte peu lorsque le processeur prédit sa cible~: le prédicteur de branchement conserve un historique des cibles,
et un chemin critique qui voit toujours le même type de message appelle toujours la même fonction. Il coûte une \omterm{def:ll:cpu-microarchitecture:branch}{erreur de prédiction de branchement} lorsque la
cible change de façon imprévisible, et il coûte toujours les optimisations qu'il bloque~: le compilateur ne peut pas inliner une fonction sans pouvoir la nommer,
ce qui l'empêche de replier le travail de la fonction appelée dans l'appelant, de garder des valeurs dans les registres à travers l'appel, ou de vectoriser
la boucle qui l'entoure. La mesure publiée du coût direct des appels virtuels en C++ (Driesen et Hölzle, 1996) le trouve
modeste en cycles~; le coût indirect, l'optimisation perdue, est en général plus élevé et n'apparaît pas dans un \omterm{def:ll:cpu-microarchitecture:ubench}{microbenchmark} de
l'appel seul.

\section{Polymorphisme statique~: templates, CRTP et concepts}

\begin{definition}[Polymorphisme statique, curiously recurring template pattern]\label{def:ll:cpp-for-latency-ii-compile-time:static}
Le \emph{polymorphisme statique}\index{polymorphisme statique} (static polymorphism) sélectionne la fonction à appeler à la compilation, par templates ou
surcharge, de sorte que chaque appel soit direct. Dans le
\emph{curiously recurring template pattern}\index{curiously recurring template pattern} (CRTP), une classe template prend sa propre classe dérivée comme paramètre, \texttt{struct Book : Base<Book>},
et la base appelle les fonctions de la classe dérivée par un \texttt{static\_cast} de \texttt{this}.
\end{definition}

Un chemin critique de trading a rarement besoin de polymorphisme à l'exécution~: l'ensemble des types de messages est fixé par le protocole, et la stratégie
est choisie au démarrage du processus. Les deux peuvent être des paramètres de template. Les concepts de C++20 énoncent l'exigence qu'un paramètre doit
satisfaire («~a \texttt{bool on(E\&)}~»), de sorte que l'erreur pour une étape inadaptée tient sur une ligne au lieu d'une page.
La \Cref{fig:ll:cpp-for-latency-ii-compile-time:dispatch} compare trois mécanismes de répartition d'un même message vers l'un de quatre gestionnaires
(ajout, annulation, exécution, transaction), chacun effectuant une ou deux additions~: gestionnaires virtuels dans une table, \texttt{std::variant} avec
\texttt{std::visit}, et CRTP avec un \texttt{switch}. Sur un flux d'un seul type de message, où chaque branchement est prédit, le
répartiteur CRTP coûte environ une demi-nanoseconde par message et le répartiteur virtuel environ quatre fois plus~; sur un mélange aléatoire de
les quatre types, les trois subissent des erreurs de prédiction, et l'écart se resserre. Les deux dernières barres sont le pipeline à quatre étapes du
début, en versions statique et dynamique.

\begin{omfigure}
\begin{tikzpicture}
\begin{axis}[width=12cm,height=5.6cm,ybar,bar width=10pt,ymin=0,
  symbolic x coords={virtual,variant,CRTP,pipeline static,pipeline dynamic},xtick=data,enlarge x limits=0.12,
  xticklabels={virtuel,variant,CRTP,{pipeline\\ statique},{pipeline\\ dynamique}},xticklabel style={align=center},
  ylabel={ns par message},tick label style={font=\small},label style={font=\small},grid=major,grid style={gray!25},
  legend columns=2,legend style={font=\footnotesize,at={(0.5,-0.33)},anchor=north,draw=none,fill=none,
  /tikz/every even column/.append style={column sep=8pt}},table/col sep=comma]
\addplot[fill=omDef!55,draw=omDef] table[x=design,y=one_type]{figdata/low-latency/07-cpp-for-latency-ii-compile-time/measured_dispatch.csv};
\addplot[fill=omThm!45,draw=omThm] table[x=design,y=mixed]{figdata/low-latency/07-cpp-for-latency-ii-compile-time/measured_dispatch.csv};
\legend{un seul type de message,mélange aléatoire de quatre}
\end{axis}
\end{tikzpicture}

\omcaption{Coût par message de trois façons de le répartir vers son gestionnaire, sur un flux d'un type et sur un mélange aléatoire, et
d'un pipeline à quatre étapes composé statiquement (\texttt{firm.pipeline}) et par appels virtuels (le coût du pipeline ne dépend pas du
mélange de messages). Mesuré sur un ordinateur portable (Intel Core Ultra 7 155H) sous WSL2, sans cœurs isolés, épinglé à un seul processeur,
meilleur de 15 exécutions de 65\,536 messages. Données~: \texttt{bench\_dispatch.py}.}\label{fig:ll:cpp-for-latency-ii-compile-time:dispatch}
\end{omfigure}

\section{Évaluation à la compilation}

\begin{definition}[Évaluation à la compilation]\label{def:ll:cpp-for-latency-ii-compile-time:consteval}
L'\emph{évaluation à la compilation}\index{evaluation@évaluation à la compilation} (compile-time evaluation) calcule une valeur pendant la compilation plutôt qu'à l'exécution~:
une fonction \texttt{constexpr} peut être évaluée par le compilateur lorsque ses arguments sont des constantes, et une fonction \texttt{consteval}
doit l'être, de sorte que son résultat soit une constante intégrée au programme.
\end{definition}

Les tables sont l'usage typique~: puissances de dix pour rééchelonner les prix entre protocoles aux décimales implicites différentes, la
correspondance entre l'octet de type de message d'un protocole et l'indice d'un gestionnaire, les longueurs des messages à taille fixe, les tables de sommes de contrôle. Construite
par le compilateur, une table ne coûte rien au démarrage, ne peut pas être oubliée dans une routine d'initialisation, et est vérifiée par
\texttt{static\_assert} à la compilation du programme. Les faits de disposition relèvent aussi de cette approche~: la construction du chapitre~3 affirme
à la compilation qu'un compteur avec remplissage occupe exactement une \omterm{def:ll:memory-hierarchy-and-caches:line}{ligne de cache}, et le chapitre~16 vérifie la taille de chaque message
par rapport à la spécification du protocole.

\section{Inlining et ses limites}

\begin{definition}[Inlining]\label{def:ll:cpp-for-latency-ii-compile-time:inline}
L'\emph{inlining}\index{inlining} est le remplacement, par le compilateur, d'un appel par le corps de la fonction appelée, ce qui supprime
le coût propre de l'appel et, plus important encore, laisse le compilateur optimiser le code combiné~: propager les constantes, garder les valeurs
dans les registres, supprimer le travail dont le résultat n'est pas utilisé.
\end{definition}

Le pipeline statique du \cref{lst:ll:cpp-for-latency-ii-compile-time:asm-static} réalise tout l'effet en treize
instructions~: les additions des quatre étapes, le ou exclusif, la comparaison et la mise à jour conditionnelle, les objets d'étape
repliés en lectures et écritures. Son jumeau dynamique (\cref{lst:ll:cpp-for-latency-ii-compile-time:asm-dyn}) est une boucle
sur un vecteur de pointeurs avec un appel indirect par étape. Le compilateur décide de l'\omterm{def:ll:cpp-for-latency-ii-compile-time:inline}{inlining} selon un modèle de coût~: petites fonctions, fonctions
appelées une seule fois, fonctions dont il voit la définition. Il n'inline pas à travers les unités de traduction sans
optimisation à l'édition de liens (LTO) (chapitre~8), à travers un pointeur de fonction qu'il ne peut résoudre, ou une fonction marquée \texttt{noinline}~; et il peut
renoncer à inliner une grande fonction dans une boucle critique. \texttt{[[gnu::always\_inline]]} outrepasse le modèle de coût pour les rares
fonctions qui doivent être inlinées~; c'est une affirmation sur les heuristiques de ce compilateur, qui doit être vérifiée dans le
code assembleur plutôt que de lui faire confiance.

\section{Disposition des branchements et chemins froids}

\begin{definition}[Chemin froid]\label{def:ll:cpp-for-latency-ii-compile-time:cold}
Un \emph{chemin froid}\index{chemin froid} (cold path) est un code qui s'exécute rarement (gestion des erreurs, journalisation d'une anomalie, reprise). Le tenir
hors de la fonction critique, en marquant la branche rare \texttt{[[unlikely]]} et la fonction rare \texttt{cold} et
\texttt{noinline}, garde le code critique petit et contigu dans le cache d'instructions, le cas courant se poursuivant sans saut.
\end{definition}

GCC place les fonctions marquées \texttt{cold} dans une section distincte avec les autres fonctions froides et traite les branchements vers elles
comme peu probables. La fonction de rééchelonnement du tutoriel vérifie ses arguments, et la vérification coûte une comparaison et une
branche prédite~; le signalement d'un mauvais argument, avec sa sortie formatée, vit hors du chemin principal. La discipline compte davantage
que n'importe quel attribut isolé~: la fonction critique doit contenir le cas courant et les appels vers tout le reste.

\section{Tutoriel~: trois répartiteurs et un pipeline}\label{tut:ll:cpp-for-latency-ii-compile-time:tut}

\begin{tutorial}
\textbf{Objectif.} Répartir les messages de trois façons, composer un pipeline à la compilation et lire dans l'assembleur ce qu'a fait le compilateur de chacune.
\textbf{État final~:} la \cref{fig:ll:cpp-for-latency-ii-compile-time:dispatch} et les deux listings d'assembleur.
\begin{tutsteps}
\item \textbf{CRTP.} La méthode \texttt{on} de la base aiguille selon le type et appelle directement le gestionnaire de la classe dérivée.
\omcode{code/low-latency/07-cpp-for-latency-ii-compile-time/cpp/ll_dispatch.hpp}{45}{64}{Répartition statique avec le curiously recurring template pattern (CRTP).}\label{lst:ll:cpp-for-latency-ii-compile-time:crtp}
\item \textbf{Une table que le compilateur construit,} et un chemin rare qu'il tient à l'écart.
\omcode{code/low-latency/07-cpp-for-latency-ii-compile-time/cpp/ll_dispatch.hpp}{66}{85}{Une table évaluée à la compilation et un chemin froid pour les erreurs.}\label{lst:ll:cpp-for-latency-ii-compile-time:consteval}
\item \textbf{Composer les étapes.} Une expression de pliage sur un tuple d'étapes, contrainte par un concept~; \texttt{\&\&} s'arrête à
la première étape qui rejette l'événement.
\omcode{code/firm/pipeline/cpp/firm_pipeline.hpp}{19}{33}{Le pipeline statique de firm.pipeline.}\label{lst:ll:cpp-for-latency-ii-compile-time:pipeline}
\item \textbf{Lire l'assembleur.} \texttt{python ll\_asm.py} compile deux fonctions d'une ligne avec
\texttt{-O2 -masm=intel} et en extrait les corps.
\omcode{code/low-latency/07-cpp-for-latency-ii-compile-time/cpp/asm/run_static.s}{1}{14}{Le pipeline statique à quatre étapes après inlining (g++ 11.4, -O2).}\label{lst:ll:cpp-for-latency-ii-compile-time:asm-static}
\omcode{code/low-latency/07-cpp-for-latency-ii-compile-time/cpp/asm/run_dynamic.s}{1}{31}{Le même pipeline par appels virtuels~: une boucle avec un appel indirect par étape.}\label{lst:ll:cpp-for-latency-ii-compile-time:asm-dyn}
\item \textbf{Mesurer} avec \texttt{python bench\_dispatch.py}.
\end{tutsteps}
\textbf{Pour aller plus loin.} Marquer les gestionnaires de la table de répartition \texttt{final} et y accéder par leurs types concrets,
et voir si le compilateur dévirtualise~; ajouter \texttt{[[gnu::noinline]]} à une étape du pipeline statique et relire
l'assembleur une nouvelle fois.
\end{tutorial}

\section{Construction~: le pipeline à la compilation}\label{bld:ll:cpp-for-latency-ii-compile-time:pipeline}

\begin{build}
\bfield{Rôle} La composition du chemin critique (hot path) de la partie~IV (décodage, carnet d'ordres, stratégie, risque, passerelle)
sous la forme d'une seule chaîne d'appels à répartition statique, avec un jumeau à polymorphisme dynamique pour les tests et pour les configurations choisies au démarrage.
\bfield{Interface} C++20 \texttt{firm::pipeline}~: concept \texttt{Stage<S, E>}~; \texttt{Pipeline<E, S\dots>} avec
\texttt{on(E\&)} et \texttt{stage<I>()}~; \texttt{AnyStage<E>}, \texttt{DynPipeline<E>} avec \texttt{add} et \texttt{on}.
Rust \texttt{firm\_pipeline}~: trait \texttt{Stage<E>}, \texttt{Chain<A, B>} et la macro \texttt{chain!},
\texttt{DynPipeline<E>}.
\bfield{Règles} Les étapes s'exécutent dans l'ordre~; le premier \texttt{false} arrête l'événement~; les pipelines statique et dynamique donnent les
mêmes résultats et le même ordre d'effets~; une étape qui ne satisfait pas le concept ne compile pas.
\bfield{Tests d'acceptation} \texttt{code/firm/pipeline/}~: ordre et court-circuit sur un événement tracé, statique contre dynamique, en
C++20 et Rust~; le test du chapitre vérifie que l'assembleur du pipeline statique ne contient aucun appel et que le pipeline dynamique
contient un appel indirect.
\bfield{Pour aller plus loin} Des étapes qui renvoient un nouveau type d'événement (un message décodé en mise à jour du carnet), avec les types de la chaîne vérifiés
par le compilateur.
\end{build}

\begin{omsources}
\item K.~Driesen et U.~Hölzle, «~The direct cost of virtual function calls in C++~», OOPSLA 1996.
\item Manuel GCC, «~Common function attributes~» (\texttt{cold}, \texttt{noinline}, \texttt{always\_inline}).
\item cppreference, \texttt{consteval}~; attributs \texttt{likely} et \texttt{unlikely}.
\item C.~Cook, «~When a microsecond is an eternity~», CppCon 2017 (modèles pour la configuration).
\end{omsources}

\section{Exercices}

\begin{exercise}[$\star$]\label{exo:ll:cpp-for-latency-ii-compile-time:1}
D'après la \cref{fig:ll:cpp-for-latency-ii-compile-time:dispatch}, combien de fois plus cher revient la \omterm{def:ll:cpp-for-latency-ii-compile-time:dyn}{répartition dynamique} que le CRTP sur un
seul type de message, et sur le mélange~?
\end{exercise}

\begin{exercise}[$\star$]\label{exo:ll:cpp-for-latency-ii-compile-time:2}
À 5 millions de messages par seconde, combien de temps de traitement par seconde le pipeline dynamique ajoute-t-il par rapport au pipeline statique~?
\end{exercise}

\begin{exercise}[$\star$]\label{exo:ll:cpp-for-latency-ii-compile-time:3}
Donner la valeur de \texttt{rescale(250, 2, 4)} et de \texttt{rescale(2500, 4, 2)}, et dire ce que fait
\texttt{rescale(1, 2, 20)}.
\end{exercise}

\begin{exercise}[$\star\star$]\label{exo:ll:cpp-for-latency-ii-compile-time:4}
Pourquoi le mélange aléatoire ralentit-il aussi le répartiteur CRTP, alors qu'il ne fait aucun appel indirect~?
\end{exercise}

\begin{exercise}[$\star\star$]\label{exo:ll:cpp-for-latency-ii-compile-time:5}
Compter les instructions du \cref{lst:ll:cpp-for-latency-ii-compile-time:asm-static} qui relèvent de chaque étape. De quelle étape
l'objet a-t-il entièrement disparu, et pourquoi~?
\end{exercise}

\begin{exercise}[$\star\star$]\label{exo:ll:cpp-for-latency-ii-compile-time:6}
Un fichier de configuration choisit au démarrage l'une de trois stratégies. Comment conserver une répartition statique~?
\end{exercise}

\begin{exercise}[$\star\star\star$]\label{exo:ll:cpp-for-latency-ii-compile-time:7}
\emph{Programmation.} Rendre une étape du pipeline statique \texttt{[[gnu::noinline]]}, régénérer l'assembleur et mesurer. Qu'est-ce qui
change dans le listing et dans le temps par événement~?
\end{exercise}

\begin{exercise}[$\star\star\star$]\label{exo:ll:cpp-for-latency-ii-compile-time:8}
\emph{Trouver la faille.} «~Un appel virtuel coûte 1.9 nanosecondes dans notre \omterm{def:ll:cpu-microarchitecture:ubench}{microbenchmark} et notre chemin en fait cinq, donc les
supprimer fait gagner au plus dix nanosecondes.~»
\end{exercise}

\section{Problème~: le branchement indirect}

\begin{problem}[{Problème du week-end --- une architecture à plug-ins sur le chemin critique}]\label{pb:ll:cpp-for-latency-ii-compile-time:1}
Un moteur de trading a été conçu avec des plug-ins~: le décodeur, le carnet d'ordres, la stratégie et le contrôle de risque sont chargés au démarrage et appelés
par une classe de base. Il traite quatre types de messages dans un mélange proche d'un tirage aléatoire. L'équipe mesure les variantes présentées dans ce
chapitre (voir la \cref{fig:ll:cpp-for-latency-ii-compile-time:dispatch}).

\medskip\noindent\textbf{Partie I --- La mesure.}
\begin{enumerate}
\item Combien coûte, par message, chacun des trois répartiteurs, sur un seul type et sur le mélange~?
\item Quel effet rend le mélange coûteux pour les trois~?
\item Combien coûte le pipeline à quatre étapes, en version statique et en version dynamique~?
\item Pourquoi le coût du pipeline ne dépend-il pas du mélange~?
\end{enumerate}

\medskip\noindent\textbf{Partie II --- Les mécanismes.}
\begin{enumerate}[resume]
\item Que fait le pipeline dynamique à chaque étape, que le statique ne fait pas (voir le \cref{lst:ll:cpp-for-latency-ii-compile-time:asm-dyn})~?
\item Combien d'instructions le pipeline statique exécute-t-il lorsque l'événement traverse toutes les étapes~?
\item Quelles optimisations sont devenues possibles une fois les étapes inlinées~?
\item Que changerait \texttt{final} sur les classes de plug-ins, et qu'est-ce qu'il ne changerait pas~?
\end{enumerate}

\medskip\noindent\textbf{Partie III --- La refonte.}
\begin{enumerate}[resume]
\item Comment des plug-ins choisis au démarrage peuvent-ils encore être répartis statiquement~?
\item Que perd-on~: qu'est-ce qui ne peut plus changer sans recompilation~?
\item Combien de variantes de programme le choix de trois stratégies et de deux contrôles de risque produit-il~?
\item Comment faut-il tester les variantes~?
\end{enumerate}

\medskip\noindent\textbf{Partie IV --- Le verdict.}
\begin{enumerate}[resume]
\item Énoncer le \emph{résultat nommé}~: le coût par message de l'appel virtuel face à la répartition statique, prévisible ou non, et le
coût du pipeline à quatre étapes dans les deux cas.
\item Pourquoi le \omterm{def:ll:cpu-microarchitecture:ubench}{microbenchmark} d'un appel virtuel seul est-il la mauvaise mesure~?
\item Quand la \omterm{def:ll:cpp-for-latency-ii-compile-time:dyn}{répartition dynamique} est-elle le bon choix dans un système de trading~?
\item Qu'apporte un concept à une interface de patrons~?
\item À quoi sert l'\omterm{def:ll:cpp-for-latency-ii-compile-time:consteval}{évaluation à la compilation} sur le chemin critique (hot path)~?
\item Pourquoi garder la gestion des erreurs hors ligne~?
\item Pourquoi vérifier l'assembleur plutôt que faire confiance à \texttt{always\_inline}~?
\item En une phrase~: que doit savoir le compilateur lorsqu'il compile le chemin critique~?
\end{enumerate}
\end{problem}

\section{Questions d'entretien}

\begin{interviewq}[$\star$ \iqroles{développeur}]\label{iq:ll:cpp-for-latency-ii-compile-time:1}
Comment fonctionne un appel de fonction virtuelle, et combien coûte-t-il~?
\end{interviewq}

\begin{interviewq}[$\star\star$ \iqroles{développeur}]\label{iq:ll:cpp-for-latency-ii-compile-time:2}
Expliquez le \omterm{def:ll:cpp-for-latency-ii-compile-time:static}{curiously recurring template pattern} et dans quels cas vous l'utiliseriez.
\end{interviewq}

\begin{interviewq}[$\star\star$ \iqroles{développeur}]\label{iq:ll:cpp-for-latency-ii-compile-time:3}
Quelle est la différence entre \texttt{constexpr} et \texttt{consteval}~?
\end{interviewq}

\begin{interviewq}[$\star\star$ \iqroles{développeur}]\label{iq:ll:cpp-for-latency-ii-compile-time:4}
Pourquoi l'\omterm{def:ll:cpp-for-latency-ii-compile-time:inline}{inlining} peut-il accélérer le code de plus que le coût de l'appel supprimé~? Dans quel cas peut-il le ralentir~?
\end{interviewq}

\begin{interviewq}[$\star\star$ \iqroles{développeur}]\label{iq:ll:cpp-for-latency-ii-compile-time:5}
\texttt{std::variant} ou une hiérarchie de classes pour les types de messages~: lequel choisiriez-vous sur un chemin critique (hot path), et pourquoi~?
\end{interviewq}

\begin{interviewq}[$\star\star\star$ \iqroles{développeur}]\label{iq:ll:cpp-for-latency-ii-compile-time:6}
Concevez un moteur de trading dont la stratégie est choisie au démarrage, mais dont le chemin critique (hot path) ne contient aucun appel indirect.
\end{interviewq}
